\documentclass{situationist-preprint}

\title{Holes, Boundaries and Hosts}
\msid{MET-15-b}
\msfamily{MET-15}
\msdiscipline{Metaphysics}
\mstype{Consequence}
\mskeywords{holes; hosts; boundaries; dependence; persistence; absences; co-location}

\begin{document}
\maketitle

\begin{abstract}
MET-15-a concludes that holes are dependent immaterial particulars located at the boundaries of their hosts. We derive five consequences concerning the relation of a hole to its host and to its boundary: a hole cannot outlive its host; it cannot change hosts; it moves with its host; it is not a part of its host; and it survives filling, which an absence cannot. The last consequence confirms that holes are not absences, so that the existence of holes is consistent with the denial of absences in MET-16.
\end{abstract}

\section{Consequences}

We take as given Corollary~3.5 of MET-15-a, that holes are dependent immaterial particulars located at the boundaries of their hosts, and Definitions~2.1 and~2.2 of that manuscript: a host is a material object whose surface has a configuration bounding a region it does not occupy, and a hole is a particular located at that region whose existence and shape depend on the configuration. MET-13-a adds that the real definition of a hole states its host and the boundary at which it is located.

\begin{corollary}[No survival without a host]\label{c:survival}
If a host ceases to exist, every hole in it ceases to exist.
\end{corollary}

\begin{proof}
A hole depends for its existence on its host (Proposition~3.4 of MET-15-a).
\end{proof}

\begin{corollary}[Host rigidity]\label{c:rigid}
No hole can be in one host at one time and in a different host at another.
\end{corollary}

\begin{proof}
It is essential to a hole to be in the host it is in, since its real definition states its host. A hole in a different host would not be that hole.
\end{proof}

Corollary~\ref{c:rigid} decides the case of division. If a slice is cut from a cheese so that a cavity lies wholly within the slice, the slice is a new host and the cavity in it is a new hole; the old hole ceased with the configuration of the old host that bounded it. Where the cheese survives the cut and the bounding surface remains part of it, the hole survives too. Cases of this kind are discussed at length by \citet{CasatiVarzi1994}; on the present account they are settled by the identity of the host.

\begin{corollary}[Co-motion]\label{c:motion}
A hole moves exactly when the portion of its host's surface that bounds it moves.
\end{corollary}

\begin{proof}
A hole is located at the region bounded by a configuration of its host's surface, and its location is fixed by that configuration.
\end{proof}

\begin{corollary}[Not a part]\label{c:part}
A hole is not a part of its host, and a host's mass, composition and material parts are what they would be if the hole were absent from the inventory.
\end{corollary}

\begin{proof}
Holes are immaterial (Proposition~3.3 of MET-15-a). Every part of a material host is made of some of its matter, and by Principle~P2 of MET-15-a whatever is made of some matter is where that matter is. A hole is where the host's matter is not.
\end{proof}

Corollary~\ref{c:part} answers the old question whether a doughnut includes its hole. It does not; it hosts it. The relation of hosting is not parthood but the dependence of a particular on the configuration of another's surface, and it is compatible with the doughnut's being a material object with only material parts. This also disposes of the worry that holes force a choice between the lining view of \citet{LewisLewis1970} and an ontology of immaterial parts of material things: holes are immaterial particulars, and they are parts of nothing material.

\begin{corollary}[Survival of filling]\label{c:fill}
A hole can be filled without ceasing to exist, and the filler and the hole are then co-located.
\end{corollary}

\begin{proof}
Filling the region a hole occupies does not change the configuration of the host's surface on which the hole depends. The filler, a material object, then occupies the region at which the hole, an immaterial particular, is located.
\end{proof}

\section{Holes are not absences}

Corollary~\ref{c:fill} distinguishes holes from absences. An absence of matter from a region ceases when matter fills the region; a hole does not. Corollaries~\ref{c:survival} and~\ref{c:rigid} distinguish them again: an absence has no host, so neither depends on one nor is rigidly tied to one. And Corollary~\ref{c:part} shows that the existence of a hole adds a particular to the world without subtracting anything from its host. The truths that MET-16 treats as truths about absences, such as that there is no milk in the refrigerator, are made true by what exists and what does not, and they require no entity. The truths about holes are positive truths about dependent particulars. The two theses are therefore consistent: there are holes and there are no absences, and the region of an unfilled hole is both the location of a hole and a region from which matter is absent, without the hole's being that absence.

\begin{thebibliography}{99}\small
\bibitem[Casati and Varzi(1994)]{CasatiVarzi1994} Casati, R. and Varzi, A.~C. (1994). \emph{Holes and Other Superficialities}. MIT Press.
\bibitem[Lewis and Lewis(1970)]{LewisLewis1970} Lewis, D. and Lewis, S. (1970). Holes. \emph{Australasian Journal of Philosophy}, 48(2), 206--212.
\end{thebibliography}

\end{document}
